% type3-words — a page of text with a Unicode mapping, and some of its words and parts of words
% in a font with none.
%
% The body is lmodern, whose fonts carry a /ToUnicode. The words in \nomap are drawn from the
% EC fonts, which under [T1] pdfTeX takes from METAFONT bitmaps and embeds as Type 3 with no
% mapping, as in type3-unmapped.tex. The extractor writes each of their glyphs [err], and the
% text around them places them, which is what glyph recovery reads them by (ADR 0048). The gold
% file holds what the source says.
\documentclass{article}
\usepackage[T1]{fontenc}
\usepackage{lmodern}
\pagestyle{empty}
\newcommand{\nomap}[1]{{\fontfamily{cmr}\selectfont #1}}
\begin{document}
Most of this page is set in a font that says what each of its glyphs is. A few
\nomap{words} are not, and neither are some \nomap{parts} of words, such as the
middle of recov\nomap{er}y or the end of glyph\nomap{s}. The \nomap{reader} has
to place each of them by the text around it.

A line can hold \nomap{two such words}, and a \nomap{number} like \nomap{1234} too.
A sign set alone, as in \nomap{\$} 5, is one glyph that nothing around it names.

The rest of the page is ordinary text, so that the page reads as one with text on it and not as
one to read whole. A page set in a font with no values throughout is read as a scan is, and its
text replaced; a page like this one keeps its text, and only the glyphs it cannot name are read
off the page. That is the case this page is for: a document whose fonts mostly say what they
draw, with a few words that do not, which is how most such documents come. The words around each
one are what place it, and the reading is taken only where they agree with it.
\end{document}
